\section{Conclusion}
\label{sec:conclusion}

For a ResNet-18 that has to recognise unseen classes from a few examples, which channels were removed mattered little. Under PMP at $30\%$ of the parameters eight scores, random selection among them, land \num{spread:pv:novel_5:PMP} to \num{spread:pd:plantdoc_5:PMP} points apart on four test settings. DACIS cannot be told from random selection, and the earlier preprint's claim for it does not carry over to held-out classes. The Taylor term was ahead on two settings, by at most \num{cmp:maxgain:PMP} points, a difference that does not survive correction for the number of comparisons. The meta stage did not raise the accuracy of the dense network, and the gain of two-stage schedules is matched by a probe that updates no weight. The numbers moved most with the class partition, the budget and the architecture: an unpruned MobileNetV3-Small beat the pruned ResNet-18 on both Earth-observation sets, and pruning bought less time (\num{cost:cpu-red:30}\% lower CPU latency) than the operation count promised (\num{cost:macs-red:30}\% fewer multiply-accumulates).

For few-shot classification of optical patches like those of RESISC45 and EuroSAT, we would start from a small pretrained architecture, and if a ResNet-18 must be pruned, use random selection or the Taylor term, skip the meta stage, and report results over several class-disjoint partitions with paired episodes. We did not prune the compact networks, try Earth-observation pretrained backbones, or measure anything on embedded hardware, so those questions remain open, as does whether a score computed on the target classes would do better.
